\documentclass[a4paper,12pt]{extarticle}

\usepackage{geometry}
\geometry{
    left=2.5cm,
    right=1.0cm,
    top=2.0cm,
    bottom=2.0cm
}

\usepackage[T2A]{fontenc}
\usepackage[utf8]{inputenc}
\usepackage[english,russian]{babel}

\usepackage{csquotes}

\usepackage{amsmath}
\usepackage{amsthm}
\usepackage{amssymb}
\usepackage{mathtools}

\usepackage{graphicx}
\graphicspath{{graphics/}}
\usepackage{float}
\usepackage{xcolor}
\usepackage{colortbl}

\usepackage{booktabs}
\usepackage{tabularx}
\usepackage{makecell}
\usepackage[flushleft]{threeparttable}
\usepackage{tablefootnote}

\usepackage{tikz}
\usepackage{pgf}
\usepackage{subcaption}

\usepackage{listings}

\usepackage{indentfirst}
\usepackage{setspace}
\usepackage{fancyhdr}

\setlength{\parindent}{1.25cm}
\onehalfspacing

\emergencystretch=3em
\tolerance=1000
\sloppy

\usepackage{chngcntr}
\counterwithin{table}{section}
\counterwithin{figure}{section}

\usepackage[
    backend=biber,
    style=numeric,
    sorting=none,
    maxbibnames=99
]{biblatex}

\addbibresource{refs.bib}

\usepackage[
    colorlinks=true,
    citecolor=blue,
    linkcolor=blue,
    urlcolor=blue,
    bookmarks=true,
    bookmarksopen=true,
    hypertexnames=true
]{hyperref}

\newcommand{\bibref}[3]{\hyperlink{#1}{#2 (#3)}}

\renewcommand{\theenumi}{\arabic{enumi}}
\renewcommand{\labelenumi}{\arabic{enumi}}
\renewcommand{\theenumii}{.\arabic{enumii}}
\renewcommand{\labelenumii}{\arabic{enumi}.\arabic{enumii}.}
\renewcommand{\theenumiii}{.\arabic{enumiii}}
\renewcommand{\labelenumiii}{\arabic{enumi}.\arabic{enumii}.\arabic{enumiii}.}





\begin{document}
    \begin{titlepage}
\newpage

{\setstretch{1.0}
\begin{center}
ПРАВИТЕЛЬСТВО РОССИЙСКОЙ ФЕДЕРАЦИИ\\
ФГАОУ ВО НАЦИОНАЛЬНЫЙ ИССЛЕДОВАТЕЛЬСКИЙ УНИВЕРСИТЕТ\\
«ВЫСШАЯ ШКОЛА ЭКОНОМИКИ»
\\
\bigskip
Факультет компьютерных наук\\
Образовательная программа «Прикладная математика и информатика»
\end{center}
}

\vspace{10em}


\begin{center}
{\bf Отчет об исследовательском проекте на тему:}\\
{\bf Исследование методов оценки перемещения БПЛА по видеопотоку при деградации или отсутствии GNSS}\\
\end{center}

\vspace{2em}

{\bf Выполнил студент: \vspace{2mm}}

{\setstretch{1.1}
\begin{tabular}{l@{\hskip 1.5cm}l}
группы \#БПМИ245, 3 курса & Злобин Герман Эдуардович \\
\end{tabular}}


\vspace{1em}
{\bf Принял руководитель проекта: \vspace{2mm}}

{\setstretch{1.1}
\begin{tabular}{l}
Басыров Эмиль Равильевич\\
Генеральный директор\\
АНО «Центр автономных решений\\
и технологий связи»
\end{tabular}}


\vspace{\fill}

\begin{center}
Москва 2027
\end{center}

\end{titlepage}


    \newpage
    \setcounter{page}{2}

    {
        \hypersetup{linkcolor=black}
        \tableofcontents
    }

    \newpage

    \section*{Аннотация}

    % TODO: написать после получения результатов

    \section*{Ключевые слова}
    БПЛА; визуальная одометрия; визуально-инерциальная навигация; SLAM; GNSS; ORB-SLAM3; VINS-Fusion; OpenVINS; MAVLink; ArduPilot.
    \pagebreak

    \section{Введение}\label{sec:introduce}

    % TODO: введение пишется после обзора и экспериментов

    \pagebreak

        \section{Обзор литературы}\label{sec:review}

    Задача оценки собственного движения по изображениям бортовой камеры изучается в робототехнике и компьютерном зрении более сорока лет.
    Для неё сложилась устойчивая терминология: \textit{визуальная одометрия} (visual odometry, VO) -- инкрементальная оценка перемещения камеры по последовательности кадров~\cite{nister_vo, vo_tutorial1}; \textit{визуально-инерциальная одометрия} (visual-inertial odometry, VIO) -- совместная обработка изображений и показаний инерциального измерительного блока (IMU)~\cite{vins_review}; \textit{SLAM} (simultaneous localization and mapping) -- одновременное построение карты среды и локализация в ней, позволяющее при повторном посещении места исправлять накопленную ошибку~\cite{slam_survey}.
    Обзоры~\cite{vo_tutorial1, vo_tutorial2, slam_survey, vins_review} систематизируют методы этих классов, работы~\cite{uav_vision_survey, avl_review} посвящены их применению на беспилотных летательных аппаратах (БПЛА).

    В данной главе сначала описано место визуальной навигации в навигационной системе БПЛА и причины, по которым одной инерциальной системы недостаточно.
    Затем изложены математические основы, общие для всех рассматриваемых методов, приведена классификация методов и обзор open-source реализаций.
    Отдельно рассмотрены особенности применения методов на БПЛА, способы передачи результата в автопилот, наборы данных и метрики оценки точности.
    Глава завершается сравнительным анализом и выбором методов для экспериментального исследования.

    \subsection{Навигация БПЛА при деградации GNSS}

    Полётный контроллер малого БПЛА оценивает положение, скорость и ориентацию аппарата расширенным фильтром Калмана (в автопилоте ArduPilot -- EKF3), объединяющим показания IMU, барометра, магнитометра и приёмника глобальной навигационной спутниковой системы (GNSS).
    IMU измеряет угловую скорость и удельную силу с частотой в сотни герц, но сам по себе даёт лишь относительную информацию: положение получается двукратным интегрированием ускорения, и любая ошибка измерения со временем накапливается.
    GNSS, напротив, даёт абсолютное положение с ошибкой порядка единиц метров, не растущей со временем, и в штатном режиме ограничивает дрейф инерциального решения.

    Сигнал GNSS слаб у поверхности Земли и легко теряется или искажается: при затенении зданиями и рельефом, многолучевом распространении, работе вблизи помех и при преднамеренном подавлении или подмене сигнала~\cite{groves}.
    Если после потери GNSS навигация опирается только на IMU, рост ошибки положения определяется в первую очередь смещениями (bias) датчиков.
    Постоянное смещение акселерометра $b_a$ даёт ошибку положения, растущую квадратично, а смещение гироскопа $b_g$ через ошибку оценки наклона и неверно скомпенсированное ускорение свободного падения $g$ -- кубически~\cite{groves}:
    \begin{equation}\label{eq:ins_drift}
        \delta p(t) \approx \frac{1}{2}\, b_a\, t^2 + \frac{1}{6}\, g\, b_g\, t^3 .
    \end{equation}
    Для MEMS-датчиков, устанавливаемых в полётные контроллеры, характерны $b_a \sim 10^{-2}$~м/с$^2$ и $b_g \sim 10$\textdegree/ч ($\approx 5\cdot 10^{-5}$~рад/с).
    Подстановка в~(\ref{eq:ins_drift}) даёт ошибку порядка $35$~м уже через минуту полёта, причём каждое слагаемое вносит около половины.
    На практике EKF частично оценивает смещения, пока доступен GNSS, но после его потери оценка смещений перестаёт уточняться, и ошибка растёт с той же скоростью.
    Поэтому без дополнительного источника информации о перемещении аппарат не может продолжать миссию дольше нескольких десятков секунд.

    Камера является естественным кандидатом на роль такого источника: она уже установлена на БПЛА, используемых для поиска людей и объектов, не излучает и не зависит от внешней инфраструктуры.
    При этом от визуального модуля не требуется заменить полётный контроллер: достаточно, чтобы он выдавал оценку относительного перемещения или скорости, которую EKF автопилота объединит с остальными датчиками.
    Именно такая постановка рассматривается в данной работе.

    \subsection{Математические основы}

    \subsubsection{Модель камеры}
    Большинство методов используют модель камеры-обскуры (pinhole).
    Точка $X_w \in \mathbb{R}^3$ в мировой системе координат проецируется в пиксель $(u, v)$ по формуле
    \begin{equation}\label{eq:pinhole}
        \lambda \begin{pmatrix} u \\ v \\ 1 \end{pmatrix} = K \left( R\, X_w + t \right), \qquad
        K = \begin{pmatrix} f_x & 0 & c_x \\ 0 & f_y & c_y \\ 0 & 0 & 1 \end{pmatrix},
    \end{equation}
    где $(R, t) \in SE(3)$ -- поза камеры, $K$ -- матрица внутренних параметров, $\lambda$ -- глубина точки~\cite{hartley_zisserman}.
    Реальные объективы вносят радиальную и тангенциальную дисторсию, которая компенсируется по параметрам, полученным при калибровке.
    Обозначим $\pi(T, X)$ функцию проекции точки $X$ камерой с позой $T = (R, t)$.

    \subsubsection{Эпиполярная геометрия и масштаб}
    Пусть $x_1, x_2$ -- нормированные координаты одной точки сцены на двух кадрах.
    Они связаны эпиполярным ограничением
    \begin{equation}\label{eq:epipolar}
        x_2^\top E\, x_1 = 0, \qquad E = [t]_\times R ,
    \end{equation}
    где $E$ -- существенная матрица, а $[t]_\times$ -- кососимметричная матрица векторного произведения.
    Матрица $E$ определена с точностью до множителя и имеет $5$ степеней свободы, поэтому относительное движение восстанавливается по пяти соответствиям~\cite{five_point}.
    Поскольку часть соответствий всегда ошибочна, оценка выполняется внутри схемы RANSAC~\cite{ransac}: многократно выбирается минимальное подмножество соответствий, строится гипотеза и подсчитывается число согласующихся с ней точек.

    Из~(\ref{eq:epipolar}) следует фундаментальное ограничение монокулярной камеры: замена $t \to s\,t$ не меняет ограничения, то есть по двум изображениям поворот определяется полностью, а перемещение -- только по направлению.
    Абсолютный масштаб траектории монокулярная система не наблюдает; более того, масштаб медленно «плывёт» вдоль траектории (scale drift)~\cite{vo_tutorial1}.
    Масштаб восстанавливается одним из способов: стереопарой с известной базой, показаниями IMU, измерением расстояния до поверхности (дальномер, барометр) или известными размерами объектов сцены.

    После того как поза первых кадров известна, положение точек сцены находится триангуляцией, а поза каждого следующего кадра -- по соответствиям «3D-точка -- пиксель» решением задачи PnP (Perspective-n-Point)~\cite{hartley_zisserman}.

    \subsubsection{Уравнивание связок}
    Точные методы уточняют позы камер $T_i$ и положения точек $X_j$ совместно, минимизируя суммарную ошибку перепроецирования (bundle adjustment)~\cite{bundle_adjustment}:
    \begin{equation}\label{eq:ba}
        \min_{\{T_i\},\, \{X_j\}} \sum_{(i,j)} \rho\!\left( \left\| z_{ij} - \pi(T_i, X_j) \right\|^2_{\Sigma_{ij}} \right),
    \end{equation}
    где $z_{ij}$ -- наблюдение точки $j$ на кадре $i$, $\Sigma_{ij}$ -- ковариация наблюдения, $\rho$ -- робастная функция потерь (например, функция Хьюбера), снижающая влияние выбросов.
    Задача~(\ref{eq:ba}) решается методами Гаусса--Ньютона или Левенберга--Марквардта с учётом разреженной структуры якобиана.
    Чтобы вычисления оставались ограниченными, в реальном времени оптимизация ведётся не по всей траектории, а по скользящему окну последних ключевых кадров; информация об исключаемых из окна переменных сохраняется в виде априорного слагаемого (маргинализация).

    \subsubsection{Модель IMU и преинтегрирование}
    Гироскоп и акселерометр измеряют угловую скорость $\omega$ и ускорение $a$ в связанной системе координат с искажениями:
    \begin{equation}\label{eq:imu}
        \tilde{\omega} = \omega + b_g + n_g, \qquad
        \tilde{a} = R_{wb}^\top \left( a_w - g_w \right) + b_a + n_a ,
    \end{equation}
    где $b_g, b_a$ -- медленно меняющиеся смещения, моделируемые случайным блужданием, $n_g, n_a$ -- белый шум, $R_{wb}$ -- ориентация аппарата, $g_w$ -- вектор ускорения свободного падения.
    IMU работает на частоте $200$--$1000$~Гц, камера -- на $20$--$30$~Гц, поэтому между соседними кадрами накапливаются десятки измерений IMU.
    Метод преинтегрирования~\cite{preintegration} сворачивает их в одно относительное ограничение на приращение ориентации, скорости и положения между ключевыми кадрами, не зависящее от абсолютной позы.
    Это позволяет включать IMU в оптимизацию~(\ref{eq:ba}) без повторного интегрирования при каждом изменении линеаризационной точки.

    \subsubsection{Наблюдаемость}
    Анализ наблюдаемости показывает, какие компоненты состояния принципиально не могут быть определены по данным датчиков.
    Для монокулярной VO ненаблюдаемы $7$ степеней свободы: глобальное положение ($3$), ориентация ($3$) и масштаб ($1$).
    Добавление IMU делает наблюдаемыми масштаб и наклон относительно вертикали (крен и тангаж), поскольку акселерометр измеряет направление гравитации; ненаблюдаемыми остаются $4$ степени свободы -- глобальное положение и рыскание~\cite{vins_consistency}.
    Следствие для данной работы: даже идеальная VIO даёт \emph{относительную} навигацию, и её ошибка по положению и курсу неизбежно накапливается.
    Ограничить этот рост можно только внешней абсолютной информацией -- возвратом GNSS, замыканием петель при повторном посещении места или сопоставлением кадров с картой местности~\cite{avl_review}.
    Ненаблюдаемые направления также требуют аккуратной реализации фильтров: при неправильной линеаризации фильтр «получает» ложную информацию о них и становится неконсистентным, то есть переоценивает свою точность~\cite{vins_consistency}.

    \subsection{Классификация методов}

    Методы визуальной и визуально-инерциальной навигации удобно классифицировать по трём независимым признакам: способу обработки изображений (фронтенд), способу оценки состояния (бэкенд) и составу датчиков.

    \subsubsection{По способу обработки изображений}

    \textbf{Методы на особых точках} (feature-based) выделяют на кадре устойчивые точки -- углы Харриса или Ши--Томаси~\cite{shi_tomasi}, FAST, ORB~\cite{orb} -- и находят их соответствия между кадрами сопоставлением дескрипторов или отслеживанием оптическим потоком Лукаса--Канаде~\cite{lucas_kanade}.
    Дальнейшая оценка движения использует только геометрию точек, то есть минимизирует ошибку перепроецирования~(\ref{eq:ba}).
    Такие методы устойчивы к изменению освещения и большим смещениям между кадрами, а дескрипторы позволяют распознавать ранее посещённые места.
    Недостаток -- зависимость от наличия выраженных углов: на однородных поверхностях точек мало.
    Классические представители -- PTAM~\cite{ptam} и семейство ORB-SLAM~\cite{orbslam, orbslam3}.

    \textbf{Прямые методы} (direct) не выделяют точки, а минимизируют фотометрическую ошибку -- разность яркости пикселей при совмещении кадров.
    Они используют всю информацию с градиентом яркости, включая края и слабую текстуру, и потому работают там, где точек мало.
    Однако прямые методы чувствительны к изменениям экспозиции, требуют точной фотометрической калибровки и хорошего начального приближения движения.
    Представители: LSD-SLAM~\cite{lsd_slam} и DSO~\cite{dso}.

    \textbf{Полупрямые методы} (semi-direct) сочетают оба подхода: начальное выравнивание выполняется по яркости небольших участков вокруг точек, а уточнение -- по геометрии.
    Представитель -- SVO~\cite{svo, svo2}, разработанный специально для БПЛА с камерой, направленной вниз, и отличающийся очень низкими вычислительными затратами.

    \subsubsection{По способу оценки состояния}

    \textbf{Фильтрационные методы} рекурсивно обновляют оценку текущего состояния и его ковариации.
    Прямое включение всех точек сцены в вектор состояния EKF ведёт к кубическому росту вычислений с числом точек.
    Фильтр MSCKF (Multi-State Constraint Kalman Filter)~\cite{msckf} решает эту проблему: в состоянии хранятся только позы нескольких последних кадров, а наблюдения точки используются как ограничение на эти позы, после чего точка исключается.
    Вычислительная сложность MSCKF линейна по числу точек, что делает его удобным для бортовых вычислителей.
    К этому классу относятся также ROVIO~\cite{rovio}, работающий с фотометрической ошибкой внутри EKF, и OpenVINS~\cite{openvins}.

    \textbf{Оптимизационные методы} решают нелинейную задачу наименьших квадратов вида~(\ref{eq:ba}) с инерциальными слагаемыми по окну ключевых кадров.
    Повторная линеаризация на каждой итерации делает их, как правило, точнее фильтров при сопоставимом размере окна, но дороже вычислительно.
    Представители: OKVIS~\cite{okvis}, VINS-Mono~\cite{vins_mono} и его развитие VINS-Fusion~\cite{vins_fusion}, Basalt~\cite{basalt}, Kimera-VIO~\cite{kimera}, ORB-SLAM3~\cite{orbslam3}.

    \textbf{Одометрия и SLAM.}
    Одометрия оценивает только локальное движение, и её ошибка растёт неограниченно.
    SLAM-системы дополнительно хранят карту ключевых кадров и выполняют \emph{замыкание петель}: при распознавании ранее посещённого места (обычно по «мешку визуальных слов» DBoW2~\cite{dbow2}) добавляется ограничение между текущей и старой позой, и вся траектория корректируется оптимизацией графа поз.
    Для задачи поиска это свойство существенно: типовые поисковые маршруты (змейка, расходящаяся спираль) многократно проходят рядом с уже обследованными участками, и замыкание петель может периодически сбрасывать накопленный дрейф.
    Цена -- больший расход памяти и вычислений, а также скачкообразные изменения оценки положения после замыкания, которые нужно корректно передавать в автопилот.

    \subsubsection{По составу датчиков}

    \textbf{Монокулярная камера} -- самый дешёвый и распространённый вариант, но масштаб в нём ненаблюдаем.

    \textbf{Стереокамера} даёт метрический масштаб за счёт известной базы между камерами.
    Однако точность определения глубины падает квадратично с расстоянием, и при высоте полёта, во много раз превышающей базу (типичная база $10$--$20$~см против высоты $30$--$100$~м), стереопара фактически вырождается в монокулярную камеру.

    \textbf{Камера и IMU} -- основной вариант для БПЛА.
    IMU обеспечивает масштаб, вертикаль и высокочастотную оценку движения при быстрых манёврах и смазанных кадрах, а камера ограничивает дрейф IMU.
    Платой является необходимость точной калибровки взаимного положения камеры и IMU и синхронизации их времени, а также инициализации: для определения масштаба и смещений требуется движение с достаточным ускорением.

    \textbf{Камера, направленная вниз, и дальномер.}
    Если поверхность под аппаратом можно считать плоской, а расстояние до неё $h$ измеряется дальномером, скорость следует непосредственно из оптического потока в центре кадра.
    В системе координат камеры для точки на оптической оси
    \begin{equation}\label{eq:flow}
        \dot{u} = -\frac{f}{h}\, v_x - f\, \omega_y
        \quad\Longrightarrow\quad
        v_x = -\frac{h}{f} \left( \dot{u} + f\, \omega_y \right),
    \end{equation}
    где $\dot{u}$ -- скорость смещения изображения в пикселях, $f$ -- фокусное расстояние в пикселях, $\omega_y$ -- угловая скорость по гироскопу; аналогично для $v_y$.
    На этом принципе построен датчик PX4FLOW~\cite{px4flow}, а автопилоты ArduPilot и PX4 штатно поддерживают такие датчики оптического потока.
    Метод не требует карты и инициализации, вычислительно очень дёшев, но измеряет только скорость, а не положение, и чувствителен к рельефу и ошибкам дальномера.
    В данной работе он рассматривается как базовый (baseline) метод, с которым сравниваются более сложные системы.

    \subsubsection{Обучаемые методы и абсолютная визуальная локализация}

    Отдельное направление составляют методы, в которых сопоставление кадров или вся оценка движения выполняются нейронной сетью.
    DROID-SLAM~\cite{droid_slam} и DPVO~\cite{dpvo} итеративно уточняют оптический поток и позы дифференцируемым уравниванием связок и на открытых наборах данных превосходят классические методы по устойчивости.
    Однако они требуют графического ускорителя, что затрудняет их использование на борту малого БПЛА, и в данной работе рассматриваются только как ориентир по точности.

    Принципиально иной подход -- абсолютная визуальная локализация~\cite{avl_review}: сопоставление кадра с заранее загруженными спутниковыми снимками или ортофотопланом местности.
    Он ограничивает ошибку, а не только замедляет её рост, но требует актуальной карты и устойчивого сопоставления разносезонных изображений.
    Этот подход находится за рамками данной работы, однако может дополнять относительную навигацию, исследуемую в ней.

    \subsection{Open-source реализации}

    Целью работы является использование готовых решений, поэтому ниже рассмотрены реализации с открытым исходным кодом, активно применяемые в исследованиях БПЛА.

    \textbf{ORB-SLAM3}~\cite{orbslam3} -- наиболее универсальная из открытых систем.
    Поддерживает монокулярную, стерео- и RGB-D-камеру, каждую как с IMU, так и без него, модели объективов pinhole и fisheye.
    Фронтенд построен на точках ORB, бэкенд -- на локальном уравнивании связок по ковизибильным ключевым кадрам.
    Инерциальная инициализация выполняется как задача оценки апостериорного максимума и занимает несколько секунд.
    Система замыкает петли и поддерживает несколько карт (Atlas): при потере трекинга начинается новая карта, которая затем объединяется со старыми при обнаружении общего места.
    По данным авторов, ORB-SLAM3 показывает одну из лучших точностей на наборе EuRoC.
    Ограничения: лицензия GPLv3, официальная поддержка только ROS~1, чувствительность к однородным текстурам из-за опоры на углы.

    \textbf{VINS-Mono и VINS-Fusion}~\cite{vins_mono, vins_fusion} -- оптимизационная VIO по скользящему окну.
    Точки выделяются детектором Ши--Томаси и отслеживаются оптическим потоком Лукаса--Канаде, что дешевле сопоставления дескрипторов.
    Инерциальные измерения учитываются через преинтегрирование, предусмотрены робастная инициализация, замыкание петель на DBoW2 и онлайн-оценка временного сдвига между камерой и IMU~\cite{vins_td}, важная для недорогих несинхронизированных датчиков.
    VINS-Fusion добавляет поддержку стереокамеры и модуль глобального объединения с GNSS, полезный для сценария, когда сигнал деградирует частично.
    Система изначально разрабатывалась для БПЛА.
    Ограничения: лицензия GPLv3, ROS~1, чувствительность к качеству калибровки.

    \textbf{OpenVINS}~\cite{openvins} -- исследовательская платформа на основе MSCKF.
    Наряду с удаляемыми точками в состояние могут включаться долгоживущие SLAM-точки.
    Поддерживается онлайн-калибровка внутренних параметров камеры, взаимного положения камеры и IMU и временного сдвига.
    Для сохранения консистентности фильтра используется техника First-Estimate Jacobians, корректно сохраняющая ненаблюдаемые направления~\cite{vins_consistency}.
    Проект хорошо документирован, поддерживает ROS~1 и ROS~2 и включает инструменты оценки траекторий и симулятор.
    Ограничения: лицензия GPLv3, отсутствие замыкания петель.

    Кроме того, при обзоре рассмотрены Basalt~\cite{basalt} (оптимизационная VIO с эффективной маргинализацией), SVO Pro~\cite{svo2} (полупрямая VO для камеры, направленной вниз), Kimera-VIO~\cite{kimera} (VIO с построением метрико-семантической карты), ROVIO~\cite{rovio} и OKVIS~\cite{okvis}.

    Независимое сравнение монокулярных VIO-систем на бортовых вычислителях БПЛА проведено в работе~\cite{vio_benchmark}.
    В ней показано, что оптимизационные системы с замыканием петель (VINS-Mono) дают наименьшую ошибку, но наиболее требовательны к ресурсам, а фильтрационные системы (MSCKF, ROVIO) уступают по точности, но значительно дешевле вычислительно.
    Общий вывод работы -- выбор метода определяется компромиссом между точностью, устойчивостью и доступными вычислительными ресурсами, а не единственной метрикой.

    \subsection{Особенности применения на БПЛА}

    Условия съёмки с БПЛА существенно отличаются от условий, в которых большинство методов разрабатывалось и тестировалось (помещения, наземный транспорт), что влияет на точность и устойчивость.

    \textbf{Высота и параллакс.}
    Чем больше высота относительно смещения между кадрами, тем меньше параллакс и тем хуже обусловлены триангуляция и оценка перемещения: при съёмке с большой высоты сцена близка к плоскости, движение камеры частично неотличимо от её поворота.
    Для монокулярных методов это ухудшает инициализацию и ускоряет дрейф масштаба.

    \textbf{Тип поверхности.}
    Над водой, снегом, ровным полем или однотонной крышей мало устойчивых точек, над водой и листвой текстура вдобавок меняется во времени.
    Над городом и пересечённой местностью точек много, но рельеф нарушает предположение о плоскости в методах на оптическом потоке.
    Зависимость ошибки от типа поверхности является одним из направлений исследования данной работы.

    \textbf{Динамика и вибрации.}
    Быстрые повороты мультикоптера и вибрация рамы вызывают смазывание кадров и, при использовании камер с построчным считыванием (rolling shutter), геометрические искажения.
    Высокочастотные вибрации также ухудшают показания IMU и могут нарушить инициализацию VIO.

    \textbf{Освещённость.}
    Тени облаков, блики и автоматическая экспозиция камеры нарушают предположение о постоянстве яркости, на котором основаны прямые методы и оптический поток.

    \textbf{Калибровка и синхронизация.}
    Для VIO необходимо знать внутренние параметры камеры, взаимное положение камеры и IMU и сдвиг их временных меток.
    Ошибка синхронизации в $10$~мс при угловой скорости $1$~рад/с эквивалентна ошибке ориентации около $0{,}6$\textdegree, что заметно ухудшает оценку.
    Поэтому важны методы с онлайн-оценкой временного сдвига~\cite{vins_td, openvins}.

    \textbf{Вычислительные ресурсы.}
    Обработка ведётся на бортовом компьютере-компаньоне (Raspberry Pi, NVIDIA Jetson и аналоги) параллельно с другими задачами, например обнаружением людей и объектов.
    Это ограничивает разрешение и частоту обработки кадров.

    \subsection{Передача навигационных данных в автопилот}

    Обмен данными между компьютером-компаньоном и полётным контроллером ведётся по протоколу MAVLink~\cite{mavlink_survey}.
    Для внешних навигационных источников в стандартном наборе сообщений предусмотрены~\cite{mavlink_common}:
    \begin{itemize}
        \item \texttt{VISION\_POSITION\_ESTIMATE} -- положение и ориентация в локальной системе координат;
        \item \texttt{VISION\_SPEED\_ESTIMATE} -- вектор скорости;
        \item \texttt{ODOMETRY} -- положение, ориентация и скорости с ковариациями, а также счётчик сбросов, позволяющий автопилоту корректно обработать скачок оценки (например, после замыкания петли);
        \item \texttt{OPTICAL\_FLOW} -- данные датчика оптического потока;
        \item \texttt{GPS\_INPUT} -- данные в формате GNSS-приёмника, позволяющие подать визуальную оценку как «виртуальный GPS».
    \end{itemize}

    В автопилоте ArduPilot источники данных для EKF3 задаются параметрами \texttt{EK3\_SRC\{1,2,3\}\_*} отдельно для горизонтального положения, вертикального положения, скорости и курса.
    Предусмотрено до трёх наборов источников с переключением между ними во время полёта, что позволяет использовать GNSS в штатном режиме и визуальную одометрию при его деградации~\cite{ardupilot_nongps, ardupilot_ekf_sources}.
    Подключение внешней системы через MAVLink включается параметром \texttt{VISO\_TYPE}.

    При интеграции требуется учитывать:
    \begin{itemize}
        \item \textbf{системы координат} -- автопилот использует NED (север--восток--вниз) и связанную систему FRD, ROS -- ENU и FLU, VO-системы -- систему координат камеры, поэтому нужно корректное преобразование с учётом положения камеры на аппарате;
        \item \textbf{задержку} -- время от съёмки кадра до получения оценки должно быть известно EKF, иначе измерение будет применено к неверному моменту;
        \item \textbf{ковариацию} -- она определяет вес, с которым EKF учитывает измерение; заниженная ковариация неконсистентного метода может ухудшить общее решение;
        \item \textbf{разрывы оценки} -- после переинициализации или замыкания петли положение меняется скачком, и автопилот должен об этом знать.
    \end{itemize}
    Для обмена сообщениями используются библиотека \texttt{pymavlink} или мост MAVROS для ROS.
    Проверка интеграции без риска для аппарата возможна в режиме программной симуляции ArduPilot SITL, где GNSS можно отключить во время полёта.

    \subsection{Наборы данных и метрики}

    \subsubsection{Наборы данных}
    Для количественной оценки требуются записи с синхронизированными изображениями, показаниями IMU и эталонной траекторией существенно более высокой точности.

    \textbf{EuRoC MAV}~\cite{euroc} -- $11$ полётов гексакоптера в машинном зале и в помещении с системой захвата движения; стереокамера $20$~Гц, IMU $200$~Гц, эталон от лазерного трекера или системы Vicon с миллиметровой точностью.
    Это основной набор для сравнения VIO-систем, результаты на нём приводятся почти во всех рассмотренных работах.
    Ограничение -- полёты в помещениях на малых высотах.

    \textbf{UZH-FPV}~\cite{uzh_fpv} -- агрессивные полёты гоночного квадрокоптера на высоких скоростях в помещении и на открытом воздухе.
    Подходит для исследования зависимости ошибки от скорости и динамики.

    \textbf{Mid-Air}~\cite{midair} -- синтетический набор полётов на малых высотах над природным ландшафтом в разные сезоны и при разной погоде, с данными камер, IMU и GNSS и точным эталоном.
    Позволяет исследовать влияние типа поверхности и условий съёмки, что сложно сделать на реальных записях.

    \textbf{KITTI}~\cite{kitti} -- стандартный набор для наземного транспорта; для БПЛА непосредственно не подходит, но широко используется для сравнения методов VO.

    \textbf{Симуляция.}
    Связка ArduPilot SITL и симулятора Gazebo позволяет порождать записи с заданными высотой, скоростью и текстурой поверхности и отключать GNSS в заданный момент, а эталонная траектория известна точно.
    Ограничение симуляции -- упрощённая модель камеры (нет смазывания, rolling shutter, шумов экспозиции), поэтому выводы требуют подтверждения на реальных данных.

    \subsubsection{Метрики}
    Оценённая траектория задаётся в собственной системе координат метода, поэтому перед сравнением её совмещают с эталонной.
    Совмещение выполняется по методу наименьших квадратов~\cite{umeyama} с числом степеней свободы, равным числу ненаблюдаемых: преобразование подобия $Sim(3)$ для монокулярной VO, сдвиг и поворот по рысканию ($4$ степени свободы) для VIO~\cite{traj_eval}.

    \textbf{Абсолютная ошибка траектории} (ATE) характеризует глобальную согласованность:
    \begin{equation}\label{eq:ate}
        \mathrm{ATE}_{\mathrm{rmse}} = \left( \frac{1}{N} \sum_{i=1}^{N} \left\| \mathrm{trans}\!\left( Q_i^{-1}\, S\, P_i \right) \right\|^2 \right)^{1/2},
    \end{equation}
    где $P_i, Q_i$ -- оценённая и эталонная позы, $S$ -- найденное совмещающее преобразование~\cite{tum_rgbd}.

    \textbf{Относительная ошибка} (RPE) оценивает дрейф на отрезках фиксированной длины $\Delta$ и не зависит от момента, в который произошла ошибка:
    \begin{equation}\label{eq:rpe}
        E_i = \left( Q_i^{-1} Q_{i+\Delta} \right)^{-1} \left( P_i^{-1} P_{i+\Delta} \right).
    \end{equation}
    Для навигационной задачи удобно выражать RPE в процентах от пройденного пути при нескольких длинах отрезков~\cite{traj_eval}.

    Для постановки задачи данной работы этих метрик недостаточно: важна не средняя ошибка по всей записи, а \emph{рост ошибки во времени после потери GNSS}.
    Поэтому дополнительно предлагается фиксировать момент отключения GNSS $t_0$, совмещать траекторию по эталону до $t_0$ (что соответствует знанию позы в момент отказа) и строить зависимость ошибки положения от $t - t_0$ и от пройденного пути.
    Для расчёта стандартных метрик используется пакет \texttt{evo}~\cite{evo}.

    \subsection{Сравнительный анализ и выбор методов}

    Характеристики рассмотренных методов сведены в таблицу~\ref{table:vo_methods}.
    Для оценки используются обозначения:
    \begin{itemize}
        \item[«$+$»] -- высокий показатель (высокая точность, низкие вычислительные затраты, хорошая поддержка);
        \item[«$+-$»] -- средний показатель;
        \item[«$-$»] -- низкий показатель или отсутствие возможности.
    \end{itemize}

    \begin{table}[H]
        \caption{Сравнение open-source методов визуальной и визуально-инерциальной навигации}
        \label{table:vo_methods}
        \scriptsize
        \centering
        \renewcommand{\arraystretch}{1.2}
        \renewcommand\tabularxcolumn[1]{m{#1}}
        \begin{tabularx}{\textwidth}{|>{\centering\arraybackslash}m{0.13\textwidth}|>{\centering\arraybackslash}m{0.11\textwidth}|>{\centering\arraybackslash}m{0.07\textwidth}|>{\centering\arraybackslash}m{0.07\textwidth}|>{\centering\arraybackslash}m{0.06\textwidth}|>{\centering\arraybackslash}X|>{\centering\arraybackslash}X|}
            \hline
            \makecell[c]{Класс/\\метод} &
            \makecell[c]{Датчики} &
            \makecell[c]{Точ-\\ность} &
            \makecell[c]{Затраты} &
            \makecell[c]{Петли} &
            \makecell[c]{Основные\\преимущества} &
            \makecell[c]{Основные\\недостатки}\\
            \hline
            \makecell[c]{Опт. поток\\\textbf{+ дальномер}} &
            моно + дальномер + гироскоп & $-$ & $+$ & $-$ &
            Минимальные затраты, штатная поддержка в автопилоте &
            Только скорость; предположение о плоской поверхности\\
            \hline
            \makecell[c]{Прямые\\\textbf{DSO}} &
            моно & $+-$ & $+-$ & $-$ &
            Работа при слабой текстуре &
            Нет масштаба и IMU; чувствительность к экспозиции\\
            \hline
            \makecell[c]{Полупрямые\\\textbf{SVO Pro}} &
            моно, стерео, + IMU & $+-$ & $+$ & $+-$ &
            Очень быстрый, ориентирован на БПЛА &
            Сложная настройка, ограниченная поддержка\\
            \hline
            \makecell[c]{Фильтр\\\textbf{ROVIO}} &
            моно + IMU & $+-$ & $+$ & $-$ &
            Низкие затраты, работа при слабой текстуре &
            Уступает по точности, проект не развивается\\
            \hline
            \makecell[c]{Фильтр\\\textbf{OpenVINS}} &
            моно, стерео, + IMU & $+$ & $+$ & $-$ &
            Онлайн-калибровка, консистентность, ROS~2, инструменты оценки &
            Нет замыкания петель\\
            \hline
            \makecell[c]{Оптимизация\\\textbf{OKVIS}} &
            стерео + IMU & $+-$ & $+-$ & $-$ &
            Строгая постановка, первая открытая оптимизационная VIO &
            Уступает новым методам\\
            \hline
            \makecell[c]{Оптимизация\\\textbf{VINS-Fusion}} &
            моно, стерео, + IMU, + GNSS & $+$ & $+-$ & $+$ &
            Создан для БПЛА; оценка сдвига времени; объединение с GNSS &
            Только ROS~1; требовательность к калибровке\\
            \hline
            \makecell[c]{Оптимизация\\\textbf{Basalt}} &
            стерео + IMU & $+$ & $+$ & $+-$ &
            Высокая точность при низких затратах, лицензия BSD &
            Требуется стереокамера\\
            \hline
            \makecell[c]{SLAM\\\textbf{ORB-SLAM3}} &
            моно, стерео, RGB-D, + IMU & $+$ & $+-$ & $+$ &
            Универсальность, мультикарты, высокая точность &
            Чувствительность к однородной текстуре; ROS~1\\
            \hline
            \makecell[c]{Обучаемые\\\textbf{DPVO}} &
            моно & $+$ & $-$ & $-$ &
            Устойчивость в сложных сценах &
            Требует GPU; нет масштаба\\
            \hline
        \end{tabularx}
    \end{table}

    \subsubsection{Выбор методов}
    Для экспериментального исследования выбраны ORB-SLAM3, VINS-Fusion и OpenVINS, а в качестве базового метода -- оценка скорости по оптическому потоку с камеры, направленной вниз, и дальномеру.

    Выбранные методы покрывают все основные классы бэкенда: ORB-SLAM3 -- SLAM на особых точках с замыканием петель и мультикартами, VINS-Fusion -- оптимизационная VIO по скользящему окну с отслеживанием точек оптическим потоком, OpenVINS -- фильтрационная VIO.
    Это позволяет сравнить не только конкретные реализации, но и подходы: оптимизацию и фильтрацию, наличие и отсутствие замыкания петель, сопоставление дескрипторов и отслеживание потоком.
    Все три системы поддерживают монокулярную камеру с IMU -- конфигурацию, наиболее реалистичную для малого БПЛА, -- что делает сравнение корректным.
    Кроме того, ORB-SLAM3 и VINS-Fusion можно запустить без IMU, что позволяет количественно оценить вклад инерциальных измерений.

    Базовый метод на оптическом потоке необходим для оценки того, оправдано ли использование сложных систем: если при типичных высотах и скоростях поискового полёта он даёт сопоставимую ошибку, его вычислительная простота и штатная поддержка в автопилоте будут решающими.
    Обучаемые методы (DPVO, DROID-SLAM) исключены из основного сравнения из-за требования к GPU, но могут использоваться как ориентир по точности на записанных данных.

    \subsection{Выводы по главе}

    Проведённый обзор показывает, что без внешней информации инерциальная навигация малого БПЛА теряет пригодность за десятки секунд, а видеопоток бортовой камеры является доступным источником данных для замедления этого роста.
    При этом ни один визуальный метод не устраняет накопление ошибки полностью: даже при использовании IMU глобальное положение и курс остаются ненаблюдаемыми, и ошибка ограничивается только замыканием петель или внешней абсолютной информацией.

    Монокулярные методы без дополнительных датчиков непригодны для навигации из-за ненаблюдаемого масштаба; стереокамера на типичных высотах полёта теряет своё преимущество.
    Наиболее подходящей конфигурацией для БПЛА является монокулярная камера с IMU, а при полёте над ровной поверхностью -- камера, направленная вниз, с дальномером.
    Точность и устойчивость методов существенно зависят от условий съёмки, характерных для БПЛА: высоты, типа поверхности, динамики полёта и качества калибровки.
    Количественных данных о такой зависимости в открытых работах мало, что и определяет исследовательскую часть данной работы.

    Для экспериментального сравнения выбраны ORB-SLAM3, VINS-Fusion и OpenVINS, представляющие разные классы методов, и базовый метод на оптическом потоке.
    Сравнение планируется проводить на наборах EuRoC и Mid-Air и в симуляции ArduPilot SITL с Gazebo по метрикам ATE, RPE и зависимости ошибки от времени после потери GNSS, а результат передавать в автопилот через сообщения MAVLink.
    \clearpage


    \printbibliography[heading=bibintoc]

\end{document}
